\documentclass[11pt]{amsart}
\usepackage[margin=1.2in]{geometry}
\usepackage{amsmath,amssymb,amsthm}
\usepackage{enumitem}
\newtheorem{lemma}{Lemma}
\newtheorem{question}{Question}
\theoremstyle{remark}
\newtheorem{remark}{Remark}
\renewcommand{\Pr}{\mathbb{P}}
\newcommand{\E}{\mathbb{E}}

\title{A fairness question for two involutions on Latin squares}
\date{October 2026}

\begin{document}
\maketitle

\section*{Setting}

Let $L$ be a Latin square of order $n$ on the symbol set $[n]$, with
rows and columns indexed by $[n]$.  Two rows, $1$ and $2$, and two
columns, $j\ne j'$, are distinguished throughout.

\paragraph{The frame.}  Let $\pi=\pi_L$ be the permutation of the
rows defined by
\[
L\bigl(\pi(r),\,j'\bigr)=L(r,j)\qquad(r\in[n]),
\]
i.e.\ $\pi(r)$ is the row in which column $j'$ holds the symbol that
column $j$ holds in row $r$.  ($\pi$ is a derangement.)  The
\emph{bit} of $L$ is
\[
\beta(L)=\mathbf 1\{\text{rows $1$ and $2$ lie in a common cycle of }\pi_L\}.
\]

\paragraph{Row-pair permutations.}  For rows $x\ne y$ let
$\rho_{x,y}$ be the permutation of the columns defined by
$L(y,\rho_{x,y}(q))=L(x,q)$.  For a cycle $Q$ of $\rho_{x,y}$ the two
rows carry the same set of symbols on the columns of $Q$.

\paragraph{Two moves.}
\begin{enumerate}[leftmargin=2.2em]
\item[(\textsc{trade})] For rows $x\ne y$ with $x,y\notin\{1,2\}$,
let $Q$ be the cycle of $\rho_{x,y}$ containing $j'$, and exchange the
entries of rows $x$ and $y$ in the columns of $Q$.
\item[(\textsc{turn})] For a cycle $c$ of $\pi_L$ containing neither
row $1$ nor row $2$, exchange the entries of columns $j$ and $j'$ in
the rows of $c$.
\end{enumerate}

\begin{lemma}
Both moves produce Latin squares, are involutions, and leave rows $1$
and $2$ unchanged.  A turn inverts $\pi$ on the cycle $c$ and leaves
its cycle partition unchanged.  A trade at $(x,y)$ replaces $\pi$ by
$(x\,y)\circ\pi$ if $j\notin Q$ (so the cycles of $x$ and $y$ are
merged if distinct and split if equal), and by
$(x\,y)\circ\pi\circ(x\,y)$ if $j\in Q$ (the cycle partition is
unchanged up to exchanging the labels $x,y$).  In particular a trade
can change $\beta$ only when $j\notin Q$; call such a pair $(x,y)$
\emph{flippable}.  Writing $b'=L(y,j')$ and $q^*$ for the column in
which row $x$ holds $b'$, the pair $(x,y)$ is flippable iff $q^*$ does
not lie on the $\rho_{x,y}$-cycle of $j$; if $\pi(x)=y$ or
$\pi(y)=x$ it is never flippable.
\end{lemma}

The proofs are direct verifications.  (Both rows carry the same
symbols on $Q$, and both columns carry the same symbols on $c$, which
gives the Latin property; the frame formulas follow from the
definition of $\pi$, since only the entries of rows $x,y$ in columns
$j,j'$ matter.)

\section*{Components and the question}

Fix a $2\times n$ Latin rectangle $R$ (the first two rows) and
columns $j\ne j'$ such that
\begin{equation}\label{eq:nonadj}
R(1,j)\ne R(2,j')\qquad\text{and}\qquad R(1,j')\ne R(2,j),
\end{equation}
i.e.\ $\pi(1)\ne2$ and $\pi(2)\ne1$ for every extension of $R$.  Let
$S(R)$ be the set of Latin squares whose first two rows are $R$.  The
two moves act on $S(R)$; let $\mathcal K(R)$ be the set of orbits
(``components'') of the group they generate, and for a component
$K$ let
\[
p(K)=\frac{|\{L\in K:\beta(L)=1\}|}{|K|}.
\]

\begin{question}[component fairness]
Is it true that, uniformly over $R$ and over column pairs $(j,j')$
satisfying \eqref{eq:nonadj},
\[
\frac1{|S(R)|}\sum_{K\in\mathcal K(R)}|K|\,\bigl(2p(K)-1\bigr)^2
\;\longrightarrow\;0\qquad(n\to\infty)?
\]
Equivalently: for $L$ uniform in $S(R)$, the conditional probability
that $\beta(L)=1$ given the component of $L$ is $\tfrac12+o(1)$ in
$L^2$.  A quantitative version asks for the left side to be
$o(1/\log^2n)$.
\end{question}

\begin{remark}[why the question is natural]
Under the uniform measure on $S(R)$ the moves define a symmetric
Markov chain (pick a uniform pair $x,y\notin\{1,2\}$ and trade, or a
uniform row $x\notin\{1,2\}$ and turn its cycle when allowed), so the
uniform measure on each component is stationary and the left side of
the Question is $\lim_{M\to\infty}\E[\beta'(L_0)\beta'(L_{2M})]$ with
$\beta'=2\beta-1$, the long-time autocorrelation of the bit.  Along
the chain the frame $\pi$ performs the walk of left multiplication by
transpositions $(x\,y)$, $x,y\notin\{1,2\}$, \emph{thinned} to the
flippable pairs (the turns and the non-flippable trades only relabel
rows within cycles).  The unthinned walk preserves $\pi^{-1}(1)$ and
$\pi^{-1}(2)$ and is the random-transposition walk on $S_{n-2}$ after
contracting the arcs $\pi^{-1}(1)\to1$ and $\pi^{-1}(2)\to2$; on each
of its cosets exactly half of the permutations have rows $1,2$ in a
common cycle.  So if every pair were flippable the answer would be
yes, with rate $o(1)$ after $O(n\log n)$ steps by Diaconis--Shahshahani.
The content of the Question is that the thinning by flippability does
not destroy this.
\end{remark}

\begin{remark}[what is known numerically]
On uniformly random Latin squares with $n=30$--$100$, a uniformly
random pair $x,y\notin\{1,2\}$ is flippable with probability
$0.50\pm0.01$, independently of the relative position of $x,y$ in the
cycles of $\pi$ except for the exclusion $\pi(x)=y$, $\pi(y)=x$; the
fraction of flippable pairs in a square is concentrated
($0.42$--$0.55$ at $n=30$); and the autocorrelation of the bit along the
chain, as a function of the number of trades performed, follows that
of the unthinned walk and decays to zero.  No square with
$|2p(K)-1|$ bounded away from $0$ has been observed except for
components of trivial size.
\end{remark}

\begin{remark}[weaker forms that are still useful]
(i) A yes in expectation over $R$ --- for $(R,j,j')$ drawn as the first
two rows and a column pair of a uniformly random Latin square, with
$(j,j')$ restricted to pairs at a prescribed distance along a cycle of
the row permutation $R(1,\cdot)\mapsto R(2,\cdot)$ --- suffices for
the application below.  (ii) A sufficient condition is that, for all
but an $o(1/\log^2n)$-fraction of $L\in S(R)$, a positive proportion
of the bit-changing pairs (the pairs $x,y$ in the common cycle of
rows $1,2$ that separate them when $\beta=1$; the pairs
$x\in C_1$, $y\in C_2$ when $\beta=0$) are flippable, together with a
mixing estimate for the thinned walk.
\end{remark}

\begin{remark}[provenance]
The Question arose in a programme on the Cavenagh--Greenhill--Wanless
conjecture on the cycle type of the second row of a random Latin
square: with $(j,j')$ a marked column pair $\{j,\sigma^\alpha(j)\}$,
$\alpha\ge2$, inside a cycle of the row permutation $\sigma$, the
marked-pair identity $2\,\mathcal C_n(\lambda)/\mathcal C_n(\mu)-1
=\Pr[\beta=1]/\Pr[\beta=0]$ reduces the conjecture to
$\Pr[\beta=1]=\tfrac12+o(1/\log n)$ over such instances, and a
positive answer to the quantitative Question, averaged as in
Remark~3(i), gives exactly that.  The question itself makes no
reference to this.
\end{remark}

\end{document}
